\subsection{DPLL SAT Solver}
% Note: Mô tả SAT/UNSAT, unit propagation, đơn giản hóa CNF, chọn biến, phân nhánh, quay lui; bổ sung giả mã và độ phức tạp đúng với cài đặt.
Sau khi các ràng buộc của các ván chơi đã được ãm hóa sang dạng chuẩn hội CNF, nhóm sử dụng thuật toán Davis--Putnam--Logemann--Loveland (DPLL) để kiểm tra công thức có nghiệm hay không. 
Trong cài đặt, mỗi literal được biểu diễn bằng một số nguyên có dấu. Chẳng hạn, \texttt{3} biểu diễn $x_3$ và \text{-3} biểu diễn $\neg x_3$. Mỗi mệnh đề là
một bộ các literal, còn công thức CNF là một bộ các mệnh đề. Hàm \texttt{solve} nhận công thức CNF và số biến, sau đó trả về một phép gán đầy đủ nếu công thức thỏa hoặc \texttt{None} nếu công thức vô nghiệm.

Ở mỗi nút tìm kiếm, bộ giải được thực hiện lan truyền đơn vị trước khi phân nhánh.
Nếu một mệnh đề chỉ còn literal $\ell$, mọi nghiệm của công thức tại nút đó đều phải làm cho $\ell$ đúng. 
Bộ giải ghi nhận phép gán này, bỏ các mệnh đề chứa $\ell$
vì chúng đã được thỏa, rồi xóa $\neg\ell$ khỏi các mệnh đề còn lại.
Quá trình được lặp lại vì một lần đơn giản hóa có thể tạo ra mệnh đề
đơn vị mới. Nếu xuất hiện một mệnh đề rỗng, nhánh hiện tại vô nghiệm.
Nếu công thức không còn mệnh đề nào, phép gán hiện tại đã thỏa toàn bộ
các ràng buộc.
Khi lan truyền đơn vị dừng mà công thức vẫn còn mệnh đề, cài đặt của
nhóm chọn biến có chỉ số nhỏ nhất còn xuất hiện trong công thức.
Bộ giải lần lượt thử giá trị \texttt{True} và \texttt{False}. Mỗi
nhánh sử dụng một bản sao của phép gán hiện tại, nhờ đó các giá trị
phát sinh ở nhánh thất bại không ảnh hưởng đến nhánh tiếp theo.
Thuật toán~\ref{alg:dpll-search} mô tả hàm tìm kiếm đệ quy.

\begin{algorithm}[htbp]
    \caption{Tìm kiếm nghiệm bằng DPLL}
    \label{alg:dpll-search}
    \begin{algorithmic}[1]
        \Require Công thức CNF $F$, phép gán hiện tại $A$
        \Ensure Một phép gán thỏa mãn hoặc $\bot$ nếu vô nghiệm
        \State $R \gets$ \Call{LanTruyenDonVi}{$F,A$}
        \If{$R=\bot$}
            \State \Return $\bot$
        \EndIf
        \State $(F,A) \gets R$
        \If{$F$ không còn mệnh đề}
            \State \Return $A$
        \EndIf
        \State $x \gets$ biến có chỉ số nhỏ nhất còn xuất hiện trong $F$
        \For{$b$ lần lượt nhận giá trị $\mathrm{True},\mathrm{False}$}
            \State $A' \gets$ bản sao của $A$ có bổ sung phép gán $x=b$
            \State $F' \gets$ công thức thu được khi đơn giản hóa $F$ theo $x=b$
            \State $R \gets$ \Call{DPLL}{$F',A'$}
            \If{$R\neq\bot$}
                \State \Return $R$
            \EndIf
        \EndFor
        \State \Return $\bot$
    \end{algorithmic}
\end{algorithm}

Xét công thức
\[
F=(\neg x_1 \lor x_2)
  \land(\neg x_1 \lor \neg x_2)
  \land(x_1 \lor x_3).
\]
Ban đầu, $F$ chưa có mệnh đề đơn vị nên bộ giải chọn $x_1$. Khi thử
$x_1=\mathrm{True}$, hai mệnh đề đầu lần lượt rút gọn thành $x_2$ và
$\neg x_2$. Hai yêu cầu này mâu thuẫn, do đó bộ giải quay lui và thử
$x_1=\mathrm{False}$. Lúc này, hai mệnh đề đầu đã được thỏa, còn mệnh
đề cuối rút gọn thành $x_3$. Lan truyền đơn vị gán
$x_3=\mathrm{True}$ và công thức không còn mệnh đề nào. Vì $x_2$
chưa cần được gán, hàm \texttt{solve} điền giá trị \texttt{False}
để trả về phép gán đầy đủ
$(x_1,x_2,x_3)=(\mathrm{False},\mathrm{False},\mathrm{True})$.

Lan truyền đơn vị bảo toàn các nghiệm vì mỗi giá trị được gán ở
bước này đều do một mệnh đề bắt buộc. Khi phải phân nhánh, hai giá trị
của biến được chọn bao phủ mọi khả năng của biến đó. Vì vậy, nếu cả
hai nhánh đều thất bại thì công thức tại nút hiện tại vô nghiệm.
Với $n$ biến, kể cả biến phụ do bộ mã hóa sinh ra, cây tìm kiếm có
thể có $O(2^n)$ nút trong trường hợp xấu nhất. Tại mỗi nút, cài đặt
còn phải quét các mệnh đề và tạo công thức rút gọn cho từng phép gán.
Lan truyền đơn vị giúp giảm số lần phân nhánh bằng cách xác định trước
những giá trị đã bị ràng buộc ép buộc.

Trong Griductive Solver, lớp \texttt{DPLLChecker} xây dựng cơ sở tri
thức $KB$ từ các manh mối đã mở và các nhận định đã được chứng minh.
Để xét một ô, hệ thống lần lượt kiểm tra $KB\land x$ và
$KB\land\neg x$, trong đó $x$ biểu diễn việc ô đó có tội. Chẳng hạn,
nếu $KB\land x$ vô nghiệm và $KB\land\neg x$ thỏa được, mọi nghiệm
phù hợp với thông tin công khai đều xác định ô đó vô tội. Kết quả
này là căn cứ để agent đưa ra gợi ý hoặc để engine chấp nhận nhận
định của người chơi.